%============================================================
%  ch11_open.tex
%  Chapter 11: Open Problems and Perspectives
%
%  Tensor Formats in Banach Spaces
%  Antonio Falcó, CEU Universities
%
%  Estimated length: ~20 pp.
%============================================================

\chapter{Open Problems and Perspectives}
\label{chap:open}

The results established in this monograph open several directions for
further research. This chapter collects the most natural open problems,
organised by theme, together with a discussion of connections to other
active areas of mathematics and its applications.

%------------------------------------------------------------
\section{Global topology of tensor manifolds}
\label{sec:open:topology}
%------------------------------------------------------------

\subsection*{Connectedness of tree-based manifolds}

We showed in Proposition~\ref{prop:TBT_connected} that each Tucker
manifold $\Tucker{\fr}{\VD}$ is connected. For tree-based manifolds
with $\dep{\TD} \geq 2$, this question is open in the
infinite-dimensional setting.

\begin{problem}
\label{prob:connected}
Let $\TD$ be a dimension partition tree with $\dep{\TD} \geq 2$ and
let $\fr \in \ADTD$ with $\FTv$ a proper set. Is $\FTv$ connected?
\end{problem}

In finite dimensions, connectedness has been established for the HT
format~\cite{UschmajewVandereycken2013} and the TT format~\cite{HoltzRohwedderSchneider2012},
but the methods used there (algebraic geometry of determinantal varieties)
do not directly extend to the Banach setting.

\subsection*{Higher homotopy groups}

Even in finite dimensions, the homotopy type of $\Tucker{\fr}{\VD}$
and $\FTv$ is not fully understood. Since $\Tucker{\fr}{\VD}$ is
a fibre bundle over the product of Grassmannians
$\prod_\alpha \Grass{r_\alpha}{V_\alpha}$ (Theorem~\ref{thm:Tucker_Banach_Manifold}),
the long exact sequence of homotopy groups relates
$\pi_n(\Tucker{\fr}{\VD})$ to $\pi_n(\prod_\alpha \Grass{r_\alpha}{\RR^{n_\alpha}})$
and $\pi_n(\RR^{\times r_\alpha}_*)$.

\begin{problem}
\label{prob:homotopy}
Compute the fundamental group $\pi_1(\Tucker{\fr}{\VD})$ and
$\pi_1(\FTv)$ in the finite-dimensional case. Are these manifolds
simply connected?
\end{problem}

\subsection*{Rank stratification}

The decomposition $\VD = \bigsqcup_\fr \Tucker{\fr}{\VD}$ is a
stratification of the algebraic tensor space. The boundary
$\overline{\Tucker{\fr}{\VD}} \setminus \Tucker{\fr}{\VD}$ consists of
tensors with strictly smaller rank, and the closure
$\overline{\Tucker{\fr}{\VD}}^{\|\cdot\|}$ in the Banach topology is the
set $\Tucker{\leq\fr}{\VD}$ of tensors of bounded rank
(Theorem~\ref{thm:weak_closed}).

\begin{problem}
\label{prob:stratification}
Describe the local geometry of the boundary strata: what is the tangent
cone to $\Tucker{\leq\fr}{\VD}$ at a boundary point
$\bv \in \Tucker{\leq\fr}{\VD} \setminus \Tucker{\fr}{\VD}$?
Is the stratification Whitney regular?
\end{problem}

%------------------------------------------------------------
\section{Riemannian and Finsler geometry on tensor manifolds}
\label{sec:open:riemannian}
%------------------------------------------------------------

In the Hilbert case, the Tucker manifold inherits a Riemannian metric
from the ambient Hilbert tensor space. Riemannian gradient methods
on $\Tucker{\fr}{\VD}$ (ALS, Riemannian gradient, retraction-based
methods) have been developed by Uschmajew and Vandereycken~\cite{UschmajewVandereycken2020}.
In the Banach setting, a natural replacement is a Finsler metric (a
continuously varying norm on each tangent space).

\begin{problem}
\label{prob:finsler}
Equip $\Tucker{\fr}{\VD}$ with the Finsler metric induced by
$\|\cdot\|_D$ (via the isomorphism
$T_\bv\Tucker{\fr}{\VD} \cong \ZD{\bv}$). Study the geodesics,
curvature, and exponential map of this Finsler manifold. In particular:
does the Dirac--Frenkel flow~\eqref{eq:DF_tucker} preserve the
Finsler structure in any sense?
\end{problem}

%------------------------------------------------------------
\section{Algorithms and convergence analysis}
\label{sec:open:algorithms}
%------------------------------------------------------------

\subsection*{Projector-splitting integrators on tree-based manifolds}

For the TT format, Lubich, Oseledets, and
Vandereycken~\cite{LubichOseledetsVandereycken2015} introduced the
\emph{projector-splitting integrator}, which integrates the
Dirac--Frenkel ODE by splitting the projection onto $\ZD{\bv_r(t)}$
into a sequence of rank-one updates. For general tree-based formats,
the corresponding integrator was introduced in~\cite{CerutiLubichWalach2020}.

The convergence analysis of these algorithms in infinite-dimensional
Banach spaces (as opposed to Hilbert spaces) requires understanding
the stability of the Dirac--Frenkel flow on $\FTv$ and the
approximation properties of the splitting.

\begin{problem}
\label{prob:integrator}
Extend the convergence analysis of the projector-splitting integrator
to the setting of Banach tensor spaces. In particular, establish
error bounds that depend explicitly on the geometry of the tree-based
manifold $\FTv$ (curvature, injectivity radius, condition number of
the rank map $\varrho_\fr$).
\end{problem}

\subsection*{Alternating least squares in Banach spaces}

The \emph{alternating least squares} (ALS) algorithm and its variants
(DMRG, HOSVD) minimise a functional over $\Tucker{\leq\fr}{\VD}$ by
alternating optimisation over one factor at a time. In Hilbert spaces,
convergence of ALS is well understood~\cite{UschmajewVandereycken2020}.
In Banach spaces, convergence requires different tools.

\begin{problem}
\label{prob:ALS}
Prove convergence of ALS for minimisation over $\Tucker{\leq\fr}{\VD}$
in a reflexive Banach tensor space $\VDnorm$. What role does the
geometry of the Tucker manifold (specifically, the structure of the
tangent spaces $\ZD{\bv}$) play in the convergence rate?
\end{problem}

%------------------------------------------------------------
\section{Beyond Banach spaces}
\label{sec:open:beyond}
%------------------------------------------------------------

\subsection*{Quasi-Banach and Fréchet spaces}

The component spaces $V_\alpha$ in this monograph are normed (Banach or
at least normed). Many function spaces of interest — $L^p$ for
$0 < p < 1$, or spaces of smooth functions — are quasi-Banach or Fréchet
but not Banach. The theory of minimal subspaces and Tucker manifolds
should extend to these settings, but the proofs require substantial
modifications.

\begin{problem}
\label{prob:quasi-banach}
Develop the theory of minimal subspaces and tensor manifolds for
quasi-Banach spaces $V_\alpha$ with $0 < p < 1$. In particular,
does the Tucker manifold $\Tucker{\fr}{\VD}$ retain its smooth
structure when the $V_\alpha$ are $L^p$ spaces for $p < 1$?
\end{problem}

\subsection*{Infinite-dimensional index sets}

All results in this monograph assume $D = \{1, \ldots, d\}$ is finite.
For problems in infinite dimensions — e.g., infinite particle quantum
systems or PDEs in infinitely many variables — one would want $D$ to
be infinite. The algebraic tensor product over an infinite index set
requires projective limits or nuclear spaces.

\begin{problem}
\label{prob:infinite_D}
Extend the tree-based tensor framework to infinite index sets
$D = \NN$ (or more general sets). What is the natural analogue of
a dimension partition tree, and what conditions on the norms guarantee
the existence of minimal subspaces and best approximations?
\end{problem}

%------------------------------------------------------------
\section{Connections with quantum information}
\label{sec:open:quantum}
%------------------------------------------------------------

\subsection*{Tensor networks and entanglement}

In quantum information, a \emph{tensor network state} on a lattice
is a tensor $\bv \in \bigotimes_{j \in D} \CC^{d_j}$ whose
transfer tensors $\{C^{(\alpha)}\}$ are specified by a graph (not
necessarily a tree). The tree-based format corresponds to
\emph{tree tensor networks} (TTN), which include MPS/TT as the
linear-chain case and MERA (multi-scale entanglement renormalisation
ansatz) as a more general tree.

The \emph{entanglement entropy} of a bipartition $D = \alpha \cup \alpha^c$
is $S(\alpha) = -\mathrm{tr}(\rho_\alpha \log \rho_\alpha)$, where
$\rho_\alpha$ is the reduced density matrix of the state on $\alpha$.
The \emph{area law} for ground states of local Hamiltonians states that
$S(\alpha)$ grows at most as $|\partial \alpha|$ (the boundary of
$\alpha$), not as $|\alpha|$ (the volume). This motivates the use of
tree-based formats: the rank constraints in $\FTv$ directly control
the entanglement structure.

\begin{problem}
\label{prob:quantum}
Formulate the Dirac--Frenkel variational principle for tree tensor
networks in the language of $C^*$-algebras and operator systems.
What is the correct generalisation of the Tucker manifold in the
non-commutative setting?
\end{problem}

%------------------------------------------------------------
\section{Learning on tensor manifolds}
\label{sec:open:learning}
%------------------------------------------------------------

\subsection*{Statistical learning theory}

In machine learning, tensor network states arise naturally as
\emph{function approximation schemes} for high-dimensional data.
The generalisation capacity of these models depends on the geometry
of the hypothesis class $\FTleq{\fr}{\VD}{\TD}$.

The \emph{Rademacher complexity} of $\FTleq{\fr}{\VD}{\TD}$ controls
the generalisation error via the Dudley entropy integral. For finite
$D$ and fixed ranks, this can be bounded using the Weyl tube formula
once the curvature of the manifold $\FTv$ is understood.

\begin{problem}
\label{prob:learning}
Compute the Rademacher complexity and metric entropy of
$\FTleq{\fr}{\VD}{\TD}$ in terms of the tree-based rank $\fr$,
the depth $\dep{\TD}$, and the dimensions $\dim V_j$. How does the
choice of tree $\TD$ affect the approximation-complexity tradeoff?
\end{problem}

\subsection*{Optimal tree selection}

Given a function $f: \prod_j I_j \to \RR$ (a target function for
approximation), which tree $\TD$ gives the best approximation of $f$
by a tree-based tensor of a given rank? This question connects the
geometry of tensor manifolds to the combinatorics of tree structures.

\begin{problem}
\label{prob:tree_selection}
Given $f \in \VDnorm$ and a bound $r$ on the tree-based rank,
find the dimension partition tree $\TD^*$ that minimises the
approximation error $\inf_{\bv \in \FTleq{\fr}{\VD}{\TD^*}} \|f - \bv\|_D$.
Is this optimisation problem tractable?
\end{problem}

